\section{Matched-Similarity Causal Test}
\label{sec:matched_similarity}

% POLISH FIX 5: Added answer-span presence analysis to interpret why HIGH-CCS
% sets hallucinate more often; the evidence points to lower answer coverage in
% HIGH-CCS contexts rather than a tested generator-confidence mechanism.
The strong CCS hypothesis predicts that, at fixed local retrieval similarity,
higher context coherence should improve faithfulness. We test that directly.
For each SQuAD query, the protocol retrieves candidate passages and chooses a
HIGH-CCS and LOW-CCS three-passage set whose mean query-passage similarity
differs by at most 0.02. The construction succeeds: mean absolute similarity
gap is 0.006 and mean CCS gap is 0.533 across 200 matched pairs.

\begin{table}[H]
\centering
\small
\caption{Matched-similarity intervention. Higher CCS does not improve
faithfulness at fixed mean query-passage similarity.}
\label{tab:causal}
\begin{tabular}{lrrrrrr}
\toprule
Setting & \(n\) & HIGH faith & LOW faith & HIGH--LOW & \(p\) & 95\% CI \\
\midrule
SQuAD / Mistral / DeBERTa & 200 & 0.636 & 0.639 & -0.002 & 0.628 & [-0.022, 0.017] \\
\bottomrule
\end{tabular}
\end{table}

The binary hallucination result points in the wrong direction for the strong
claim: HIGH-CCS contexts hallucinate more often than LOW-CCS contexts
(16.5\% vs.\ 9.0\%). We therefore reject the causal CCS interpretation tested
here. CCS may still be diagnostically useful, but it is not sufficient to
identify answer support once local retrieval similarity is held fixed.

A span-presence check suggests a simple reason for the wrong-sign result. We
normalize the gold answer and ask whether it appears anywhere in the three
retrieved passages. LOW-CCS contexts contain the answer more often than
HIGH-CCS contexts (27.0\% vs.\ 17.5\%; exact McNemar \(p=0.011\)); the
four-way count is 19 pairs with the answer in both, 35 in LOW only, 16 in HIGH
only, and 130 in neither.

Thus the HIGH-CCS sets do not merely provide an equally answer-bearing but
more coherent scaffold; they more often omit the answer span. A plausible
reading is that coherent-looking context can still be incomplete, encouraging
fluent unsupported answers when the needed span is absent
\citep{maynez2020faithfulness}. We report this as a hypothesis suggested by
the data, not as a tested mechanism.

\begin{figure}[H]
\centering
\includegraphics[width=0.75\linewidth]{fix_01_paired_diff.pdf}
\caption{Paired HIGH--LOW CCS faithfulness differences under the
matched-similarity intervention. The diagonal line represents equal HIGH and
LOW faithfulness; points above the line are pairs where HIGH-CCS helped, and
points below it are pairs where it hurt. The distribution straddles the
diagonal, so the controlled test rejects the strong CCS hypothesis rather than
supporting it.}
\label{fig:matched_ccs_null}
\end{figure}
